%% Supplementary Material — Benchmark Model Configurations
%% Compile with: pdflatex supplementary_benchmark.tex
%%
%% Required packages: booktabs, tabularx, multirow, xcolor, geometry,
%%                    hyperref, array, makecell, caption, parskip

\documentclass[10pt,a4paper]{article}

\usepackage[margin=2.5cm]{geometry}
\usepackage{booktabs}
\usepackage{tabularx}
\usepackage{multirow}
\usepackage{array}
\usepackage{makecell}
\usepackage{caption}
\usepackage{xcolor}
\usepackage{hyperref}
\usepackage{parskip}
\usepackage{microtype}
\usepackage{lmodern}
\usepackage[T1]{fontenc}

\hypersetup{colorlinks=true, linkcolor=blue, urlcolor=blue, citecolor=blue}

\newcolumntype{L}[1]{>{\raggedright\arraybackslash}p{#1}}
\newcolumntype{C}[1]{>{\centering\arraybackslash}p{#1}}

\setlength{\tabcolsep}{5pt}
\renewcommand{\arraystretch}{1.25}

% ── convenience command for monospace HF IDs ──────────────────────────
\newcommand{\hfid}[1]{\texttt{\small #1}}
\newcommand{\na}{---}

\title{\textbf{Supplementary Material}\\[4pt]
       \large Benchmark Model Configurations for\\
       Shortcut Learning in Radiology Report Generation}
\date{}

\begin{document}
\maketitle
\tableofcontents
\clearpage

% ═══════════════════════════════════════════════════════════════════════
\section{Benchmark Overview}
\label{sec:overview}
% ═══════════════════════════════════════════════════════════════════════

We evaluate nine vision-language models (VLMs) on the PadChest-GR dataset
using a unified inference pipeline built on HuggingFace Transformers
4.40$+$ and PyTorch 2.1$+$.
All experiments are run on the NAISS Alvis HPC cluster (NAISS project
\texttt{NAISS2025-5-662}) with A40/A100 GPUs via SLURM.
A fixed random seed of 3 is applied throughout (\texttt{numpy},
\texttt{random}, and \texttt{torch}).

\paragraph{Occlusion experiments.}
Each model is evaluated under five conditions:
\begin{itemize}
  \item \textbf{Baseline} — full, unmodified image.
  \item \textbf{OCO-$p$} (Object-Class Occlusion) — the bounding-box region
        of the target pathology is occluded at ratio $p \in \{20, 40, 60,
        80, 100\}\%$.
  \item \textbf{DOCO-$p$} (Dataset-level OCO) — occlusion of a co-occurring
        pathology region at the same ratios.
  \item \textbf{RO-$p$} (Random Occlusion) — a randomly positioned patch of
        the same size, used as a noise-floor reference.
\end{itemize}

\paragraph{Evaluation metrics.}
Reports are scored with BLEU-1/2/4, ROUGE-L, RadGraph-F1 (simple, partial,
complete), and CheXBERT classification F1 per disease class, as implemented
in \texttt{evaluations/run\_eval.py}.

% ═══════════════════════════════════════════════════════════════════════
\section{Model Overview}
\label{sec:models}
% ═══════════════════════════════════════════════════════════════════════

Table~\ref{tab:model_overview} lists the nine evaluated models with their
HuggingFace Hub identifiers, numerical precision, and the Python virtual
environment used for inference.
Because of irreconcilable dependency conflicts between model families, five
separate environments are maintained (see Table~\ref{tab:venvs}).

\begin{table}[ht]
\centering
\caption{Model overview. HF ID = HuggingFace Hub model identifier;
         dtype = floating-point precision during inference.}
\label{tab:model_overview}
\small
\begin{tabular}{L{2.4cm} L{6.5cm} C{1.6cm} C{2.2cm}}
\toprule
\textbf{Model} & \textbf{HuggingFace Hub ID} & \textbf{dtype} & \textbf{Environment} \\
\midrule
MedGemma       & \hfid{google/medgemma-1.5-4b-it}                                             & bfloat16 & \texttt{.venv\_RRG}        \\
MAIRA-2        & \hfid{microsoft/maira-2}                                                      & bfloat16 & \texttt{.venv\_RRG}        \\
CXRMate        & \hfid{aehrc/cxrmate-rrg24}                                                    & float32  & \texttt{.venv\_RRG}        \\
NV-Reason-CXR  & \hfid{nvidia/NV-Reason-CXR-3B}                                                & bfloat16 & \texttt{.venv\_nv}         \\
CheXOne        & \hfid{StanfordAIMI/CheXOne}                                                   & bfloat16 & \texttt{.venv\_nv}         \\
LIBRA-v1       & \hfid{X-iZhang/libra-v1.0-7b}                                                & bfloat16 & \texttt{.SC\_Libra\_venv}  \\
LIBRA-LLaVA    & \hfid{X-iZhang/libra-llava-rad}                                               & bfloat16 & \texttt{.SC\_Libra\_venv}  \\
CheXagent-2    & \hfid{StanfordAIMI/CheXagent-2-3b-srrg-findings}                              & bfloat16 & \texttt{.venv\_chexagent}  \\
RaDialog       & \hfid{Chantal/RaDialog-interactive-radiology-report-generation}               & bfloat16 & \texttt{.radialog\_venv}   \\
\bottomrule
\end{tabular}
\end{table}

\begin{table}[ht]
\centering
\caption{Python virtual environments and corresponding requirement files.}
\label{tab:venvs}
\small
\begin{tabular}{L{3.0cm} L{4.5cm} L{4.5cm}}
\toprule
\textbf{Environment} & \textbf{Requirements file} & \textbf{Models} \\
\midrule
\texttt{.venv\_RRG}       & \texttt{requirements\_rrg.txt}         & MedGemma, MAIRA-2, CXRMate \\
\texttt{.venv\_nv}        & \texttt{requirements\_reasoning.txt}   & NV-Reason-CXR, CheXOne \\
\texttt{.SC\_Libra\_venv} & \texttt{requirements\_llava-libra.txt} & LIBRA-v1, LIBRA-LLaVA \\
\texttt{.venv\_chexagent} & \texttt{requirements\_chexagent.txt}   & CheXagent-2 \\
\texttt{.radialog\_venv}  & \texttt{requirements\_radialog.txt}    & RaDialog \\
\bottomrule
\end{tabular}
\end{table}

% ═══════════════════════════════════════════════════════════════════════
\section{Decoding Hyperparameters}
\label{sec:decoding}
% ═══════════════════════════════════════════════════════════════════════

Table~\ref{tab:decoding} reports the text-generation hyperparameters
used for each model.
Where a parameter is not applicable (e.g., temperature under greedy
decoding), a dash (\na) is shown.

\begin{table}[ht]
\centering
\caption{Decoding hyperparameters per model. Strategy: G = greedy search,
         BS = beam search, S = sampling.}
\label{tab:decoding}
\small
\begin{tabular}{L{2.4cm} C{1.5cm} C{1.8cm} C{1.6cm} C{1.4cm} C{1.4cm}}
\toprule
\textbf{Model} &
\textbf{Strategy} &
\textbf{max\_new\_tokens} &
\textbf{num\_beams} &
\textbf{temp.} &
\textbf{top\_p} \\
\midrule
MedGemma      & G  & 450   & 1 & \na  & \na  \\
MAIRA-2       & G  & 450   & 1 & \na  & \na  \\
CXRMate$^{a}$ & BS & \na   & 4 & \na  & \na  \\
NV-Reason-CXR & G  & 2048  & 1 & \na  & \na  \\
CheXOne       & G  & 1024  & 1 & \na  & \na  \\
LIBRA-v1      & S  & 300   & 1 & 0.20 & \na  \\
LIBRA-LLaVA   & S  & 1024  & 5 & 0.20 & \na  \\
CheXagent-2   & G  & 512   & 1 & \na  & \na  \\
RaDialog      & G  & 300   & 1 & \na  & \na  \\
\bottomrule
\end{tabular}
\vspace{4pt}\\
\raggedright\small
$^{a}$ CXRMate uses beam search with \texttt{bad\_words\_ids} to suppress
non-findings tokens (\texttt{[NF]}, \texttt{[NI]}).
Output is post-processed via \texttt{split\_and\_decode\_sections()} to
extract the findings section.
\end{table}

% ═══════════════════════════════════════════════════════════════════════
\section{Inference Prompts}
\label{sec:prompts}
% ═══════════════════════════════════════════════════════════════════════

Table~\ref{tab:prompts} lists the exact prompt string passed to each model
(key \texttt{<model>\_default} in \texttt{src/benchmark/prompts.py}).
An empty entry indicates that the model uses no free-text prompt; instead,
the image is passed directly to the model's native interface.

\begin{table}[ht]
\centering
\caption{Inference prompt per model (verbatim from \texttt{prompts.py}).}
\label{tab:prompts}
\small
\begin{tabular}{L{2.4cm} L{11.0cm}}
\toprule
\textbf{Model} & \textbf{Prompt} \\
\midrule
MedGemma
  & \textit{``Describe this X-ray.''} \\
\midrule
MAIRA-2
  & \textit{(none — structured API: frontal image $+$ optional indication /
    technique / comparison fields)} \\
\midrule
CXRMate
  & \textit{(none — encoder-decoder; report generated directly from image
    features)} \\
\midrule
NV-Reason-CXR
  & \textit{``Find abnormalities and support devices.''} \\
\midrule
CheXOne
  & \textit{``Write an example findings section for the CXR.''} \\
\midrule
LIBRA-v1
  & \textit{``Provide a detailed description of the findings in the
    radiology image.''} \\
\midrule
LIBRA-LLaVA
  & \textit{``Describe the findings in the radiology image.''} \\
\midrule
CheXagent-2
  & \textit{``Structured Radiology Report Generation for Findings
    Section''} \\
\midrule
RaDialog
  & \textit{``You are to act as a radiologist and write the finding section
    of a chest x-ray radiology report for this X-ray image.
    Write in the style of a radiologist, write one fluent text without
    enumeration, be concise and don't provide explanations or reasons.''} \\
\bottomrule
\end{tabular}
\end{table}

% ═══════════════════════════════════════════════════════════════════════
\section{Image Preprocessing}
\label{sec:preprocessing}
% ═══════════════════════════════════════════════════════════════════════

Table~\ref{tab:preprocessing} summarises the image preprocessing pipeline
applied before inference.
Most models receive a three-channel uint8 RGB image normalised with
CXR-specific min-max normalisation (implemented in
\texttt{src/benchmark/preprocess.py}).
Exceptions are detailed below.

\begin{table}[ht]
\centering
\caption{Per-model image preprocessing.}
\label{tab:preprocessing}
\small
\begin{tabular}{L{2.4cm} L{11.2cm}}
\toprule
\textbf{Model} & \textbf{Preprocessing pipeline} \\
\midrule
MedGemma
  & CXR min-max normalisation $\to$ uint8 RGB; standard HF processor. \\
\midrule
MAIRA-2
  & CXR min-max normalisation $\to$ uint8 RGB;
    \texttt{format\_and\_preprocess\_reporting\_input()} from the MAIRA-2
    processor (frontal only; no lateral image provided). \\
\midrule
CXRMate
  & \texttt{torchvision.transforms.v2}: PILToTensor $\to$ Grayscale (3-ch)
    $\to$ Resize to \texttt{encoder.image\_size} $\to$ CenterCrop $\to$
    Normalize(\texttt{encoder.image\_mean}, \texttt{encoder.image\_std}).
    Findings and impression decoded separately. \\
\midrule
NV-Reason-CXR
  & CXR min-max normalisation $\to$ uint8 RGB;
    \texttt{apply\_chat\_template()} with image token. \\
\midrule
CheXOne
  & CXR min-max normalisation $\to$ uint8 RGB;
    Qwen2.5-VL processor with left-padded tokeniser. \\
\midrule
LIBRA-v1
  & Libra \texttt{load\_pretrained\_model()} image processor;
    two-slot TAC projector (current image only; prior slot left blank). \\
\midrule
LIBRA-LLaVA
  & Libra \texttt{load\_pretrained\_model()} image processor;
    two-slot input: \texttt{[current, prior]}; conversation template
    \texttt{v1}; sampling with stopping-criteria at separator tokens. \\
\midrule
CheXagent-2
  & CXR min-max normalisation $\to$ uint8 RGB;
    images written to temp files and loaded via
    \texttt{tokenizer.from\_list\_format()};
    system prompt: \textit{``You are a helpful assistant.''}\\
\midrule
RaDialog
  & Resize to 512 px $\to$ CenterCrop $488{\times}488$ $\to$
    ToTensor $\to$ ExpandChannels (grayscale $\to$ 3-ch);
    CheXpert classifier run first to predict findings labels, which
    are prepended to the text prompt; LLaVA \texttt{vicuna\_v1}
    conversation template. \\
\bottomrule
\end{tabular}
\end{table}

% ═══════════════════════════════════════════════════════════════════════
\section{Common Experimental Settings}
\label{sec:common}
% ═══════════════════════════════════════════════════════════════════════

\begin{table}[ht]
\centering
\caption{Common settings shared across all model experiments.}
\label{tab:common}
\small
\begin{tabular}{L{4.5cm} L{9.0cm}}
\toprule
\textbf{Setting} & \textbf{Value} \\
\midrule
Dataset            & PadChest-GR \\
Random seed        & 3 (\texttt{numpy}, \texttt{random}, \texttt{torch}) \\
Hardware           & NVIDIA A40 / A100 (NAISS Alvis cluster) \\
Scheduler          & SLURM (project \texttt{NAISS2025-5-662}) \\
Python version     & 3.11.5 (\texttt{Python/3.11.5-GCCcore-13.2.0}) \\
Framework          & PyTorch 2.1$+$, HuggingFace Transformers 4.40$+$ \\
Default precision  & bfloat16 (except CXRMate: float32) \\
Output format      & JSONL per model/experiment \\
Occlusion levels   & $p \in \{0, 20, 40, 60, 80, 100\}\%$ \\
Evaluation metrics & BLEU-1/2/4, ROUGE-L, RadGraph-F1,
                     CheXBERT per-class F1 \\
\bottomrule
\end{tabular}
\end{table}

\end{document}
